\chapter{1921:Frederick Soddy}
\label{chap:chemistry1921}

The Nobel Prize in Chemistry for the year 1921 was awarded to Frederick Soddy.

\textbf{Motivation:}
for his contributions to our knowledge of the chemistry of radioactive substances, and his investigations into the origin and nature of isotopes

\section{Frederick Soddy}

\subsection*{Early Life and Education}

Frederick Soddy was born on 1877-09-02 in Eastbourne, United Kingdom.
Frederick Soddy FRS was an English radiochemist who explained, with Ernest Rutherford, that radioactivity is due to the transmutation of elements, now known to involve nuclear reactions. He also proved the existence of isotopes of certain radioactive elements. In 1921, he received the Nobel Prize in Chemistry "for his contributions to our knowledge of the chemistry of radioactive substances, and his investigations into the origin and nature of isotopes". Soddy was a polymath who mastered chemistry, nuclear physics, statistical mechanics, finance, and economics.

\subsection*{Career and Major Research}

Soddy worked at McGill with Rutherford, where the two proposed the disintegration theory of radioactivity, and later at Glasgow and Oxford. He showed that radioactive decay involves a change of element, and he established that elements can occur in several forms with different atomic weights but identical chemical properties, which he named isotopes.

\subsection*{Nobel Prize-winning Work}

The work for which Frederick Soddy was awarded the Nobel Prize in Chemistry in 1921 was described by the Nobel Committee as: "for his contributions to our knowledge of the chemistry of radioactive substances, and his investigations into the origin and nature of isotopes".

This research fundamentally advanced the field by making groundbreaking contributions that transformed chemical science.

\subsection*{Significance and Impact}

The impact of Frederick Soddy's work extends far beyond the specific achievement recognized by the Nobel Prize.
These contributions continue to influence chemical research and education today.

\subsection*{Later Career and Honors}

Frederick Soddy continued his scientific work until 1956-09-22, passing away in Brighton, United Kingdom.
He received numerous honors including membership in major scientific academies and societies.

\subsection*{Legacy}

The legacy of Frederick Soddy endures through the lasting impact of his scientific contributions.
Frederick Soddy's work continues to be cited and built upon by researchers across the chemical sciences.

\subsection*{Modern Developments}

Soddy's concept of isotopes became central to modern chemistry and physics. Radioactive isotopes are now routine tools in nuclear medicine, medical imaging, and cancer therapy, while stable isotope techniques trace metabolic pathways and date geological and archaeological materials. The discovery of isotopes also made possible nuclear energy and the modern understanding of atomic structure.

\subsection*{Selected Publications}

Frederick Soddy's major publications include numerous papers in leading journals of the era, detailing the experimental and theoretical work summarized above.

\clearpage

\section*{Chapter Summary}

Frederick Soddy was honored for his contributions to the chemistry of radioactive substances and his investigations into the origin and nature of isotopes.
